\documentclass[a4paper,landscape,8pt]{extarticle}
\usepackage[landscape,margin=0.7cm,top=0.5cm,bottom=0.7cm]{geometry}
\usepackage[utf8]{inputenc}
\usepackage[T1]{fontenc}
\usepackage[ngerman]{babel}
\usepackage{helvet}
\renewcommand{\familydefault}{\sfdefault}
\usepackage{xcolor}
\usepackage{tikz}
\usetikzlibrary{positioning, shapes.geometric}
\usepackage{enumitem}
\usepackage{multicol}
\usepackage{array}
\usepackage{fancyhdr}
\usepackage{amssymb}

% Farben
\definecolor{safariblue}{HTML}{007AFF}
\definecolor{addressgreen}{HTML}{34C759}
\definecolor{tabsorange}{HTML}{FF9500}
\definecolor{bookmarkpurple}{HTML}{AF52DE}
\definecolor{lightgray}{HTML}{F5F5F5}
\definecolor{bordergray}{HTML}{BBBBBB}
\definecolor{warnred}{HTML}{C62828}
\definecolor{privategray}{HTML}{5856D6}

\pagestyle{fancy}
\fancyhf{}
\renewcommand{\headrulewidth}{0pt}
\fancyfoot[C]{\footnotesize Seite \thepage{} von 3}

\setlength{\parindent}{0pt}
\setlength{\parskip}{2pt}

% Kompakte Listen
\setlist{nosep, leftmargin=1em, itemsep=0pt, topsep=1pt}

\begin{document}

% ============================================================================
% SEITE 1: INFORMATIONSBLATT
% ============================================================================

\noindent{\Large\bfseries Browser \& Internet: Mit Safari im Web surfen} \hfill
{\footnotesize\textbf{Lernziele:} Webseiten aufrufen | Lesezeichen nutzen | Mit Tabs arbeiten}

\vspace{1mm}
\hrule
\vspace{1mm}

% Drei-Spalten-Layout
\begin{minipage}[t]{0.32\linewidth}
\begin{center}
\colorbox{safariblue!10}{%
\begin{minipage}{0.95\linewidth}
\vspace{2mm}
\begin{center}
{\Large\bfseries\textcolor{safariblue}{Was ist ein Browser?}}\\[1pt]
{\footnotesize\itshape Dein Fenster ins Internet}
\end{center}
\vspace{2mm}
\end{minipage}%
}
\end{center}

\vspace{1mm}

\textbf{Was ist ein Browser?} Ein \textbf{Programm zum Anzeigen} von Webseiten im Internet.

\textbf{Analogie:} Wie ein \textbf{Fernseher} -- du brauchst ihn, um die Sender (Webseiten) zu sehen.

\textbf{Safari:} Der Browser von Apple -- auf jedem iPhone, iPad und Mac vorinstalliert.

\textbf{Symbol:} Ein blauer Kompass

\vspace{2mm}

\textbf{Andere Browser:}
\begin{itemize}
\item \textbf{Google Chrome} -- von Google
\item \textbf{Firefox} -- kostenlos, datenschutzfreundlich
\item \textbf{Edge} -- von Microsoft
\end{itemize}

Alle können das Gleiche -- Safari ist auf Apple-Geräten am besten integriert.

\vspace{2mm}

\textbf{Safari öffnen:}
\begin{itemize}
\item Tippe auf das \textbf{Safari-Symbol}
\item Oder: Spotlight-Suche → \glqq Safari\grqq{}
\end{itemize}

\end{minipage}%
\hfill%
\begin{minipage}[t]{0.32\linewidth}
\begin{center}
\colorbox{addressgreen!10}{%
\begin{minipage}{0.95\linewidth}
\vspace{2mm}
\begin{center}
{\Large\bfseries\textcolor{addressgreen}{Webseiten aufrufen}}\\[1pt]
{\footnotesize\itshape Adressen eingeben und surfen}
\end{center}
\vspace{2mm}
\end{minipage}%
}
\end{center}

\vspace{1mm}

\textbf{Die Adressleiste:} Das \textbf{Eingabefeld oben} -- hier gibst du Webadressen ein.

\textbf{Analogie:} Wie die \textbf{Hausnummer} -- jede Webseite hat eine eindeutige Adresse.

\vspace{2mm}

\textbf{Webseite aufrufen:}
\begin{enumerate}
\item Tippe in die \textbf{Adressleiste}
\item Gib die Adresse ein (z.B. \texttt{www.tagesschau.de})
\item Tippe auf \textbf{\glqq Öffnen\grqq{}} auf der Tastatur
\end{enumerate}

\textbf{Tipp:} Du musst \texttt{https://} nicht eingeben -- Safari ergänzt das automatisch.

\vspace{2mm}

\textbf{Auf der Seite navigieren:}
\begin{itemize}
\item \textbf{Wischen} = Scrollen (hoch/runter)
\item \textbf{Tippen auf Link} = Neue Seite öffnen
\item \textbf{Zwei Finger auseinander} = Zoomen
\end{itemize}

\textbf{Zurück / Vor:}
\begin{itemize}
\item \textbf{Pfeil links} = Zurück zur vorigen Seite
\item \textbf{Pfeil rechts} = Wieder vorwärts
\item Oder: \textbf{Vom Rand wischen}
\end{itemize}

\end{minipage}%
\hfill%
\begin{minipage}[t]{0.32\linewidth}
\begin{center}
\colorbox{lightgray}{%
\begin{minipage}{0.95\linewidth}
\vspace{2mm}
\begin{center}
{\Large\bfseries Safari-Oberfläche}\\[1pt]
{\footnotesize\itshape Die wichtigsten Elemente}
\end{center}
\vspace{2mm}
\end{minipage}%
}
\end{center}

\vspace{1mm}

\textbf{Am iPhone (von oben nach unten):}

\renewcommand{\arraystretch}{1.3}
\begin{tabular}{@{}p{2.5cm}p{4.5cm}@{}}
\textbf{Adressleiste} & Oben (oder unten bei neueren iOS) \\
\textbf{Teilen-Button} & Quadrat mit Pfeil nach oben \\
\textbf{Lesezeichen} & Buch-Symbol \\
\textbf{Tabs} & Zwei überlappende Quadrate \\
\textbf{Neue Seite} & + Symbol \\
\end{tabular}
\renewcommand{\arraystretch}{1}

\vspace{3mm}

\textbf{Adressleiste unten?}\\
Bei neueren iPhones ist die Adressleiste \textbf{unten} -- leichter zu erreichen!

\textbf{Umschalten:}\\
Einstellungen → Safari → \glqq Tableiste\grqq{} (Oben/Unten)

\vspace{2mm}

\begin{center}
\fcolorbox{safariblue}{safariblue!8}{%
\begin{minipage}{0.92\linewidth}
\vspace{2mm}
\textbf{Tipp: Seite aktualisieren}\\[2pt]
Von oben nach unten ziehen (wie bei E-Mail) -- die Seite wird neu geladen.
\vspace{2mm}
\end{minipage}}
\end{center}

\end{minipage}

\vspace{1mm}

% Zusammenfassung
\begin{center}
\fcolorbox{bordergray}{lightgray}{%
\begin{minipage}{0.98\linewidth}
\centering
\vspace{1.5mm}
\textbf{Merke:} Safari ist dein \textbf{Tor ins Internet}. Adresse eingeben → Seite ansehen → Links folgen.
\vspace{1.5mm}
\end{minipage}%
}
\end{center}

\newpage

% ============================================================================
% SEITE 2: LESEZEICHEN & TABS
% ============================================================================

\begin{center}
{\LARGE\bfseries Lesezeichen, Tabs \& Verlauf}\\[3pt]
{\normalsize Webseiten speichern und organisieren}
\end{center}

\vspace{1mm}
\hrule
\vspace{2mm}

\begin{multicols}{2}

\section*{\textcolor{bookmarkpurple}{Lesezeichen -- Deine Favoriten}}

\textbf{Was sind Lesezeichen?} Gespeicherte \textbf{Links zu Webseiten}, die du oft besuchst.

\textbf{Analogie:} Wie \textbf{Eselsohren} im Buch -- markieren wichtige Stellen.

\vspace{2mm}

\textbf{Lesezeichen speichern:}
\begin{enumerate}
\item Öffne die gewünschte Webseite
\item Tippe auf das \textbf{Teilen-Symbol} (Quadrat mit Pfeil)
\item Wähle \textbf{\glqq Lesezeichen hinzufügen\grqq{}}
\item Namen anpassen (optional)
\item \textbf{\glqq Sichern\grqq{}} tippen
\end{enumerate}

\textbf{Lesezeichen öffnen:}
\begin{enumerate}
\item Tippe auf das \textbf{Buch-Symbol} (unten)
\item Wähle \textbf{Lesezeichen} (oben)
\item Tippe auf das gewünschte Lesezeichen
\end{enumerate}

\textbf{Lesezeichen organisieren:}
\begin{itemize}
\item \textbf{Ordner erstellen} für Themen (Nachrichten, Shopping...)
\item \textbf{Bearbeiten} tippen zum Löschen/Sortieren
\end{itemize}

\vspace{2mm}

\textbf{Favoriten:} Besondere Lesezeichen, die in der \textbf{Startseite} angezeigt werden -- noch schneller erreichbar.

\vspace{3mm}

\section*{\textcolor{tabsorange}{Tabs -- Mehrere Seiten gleichzeitig}}

\textbf{Was sind Tabs?} \textbf{Mehrere Webseiten} gleichzeitig geöffnet -- wie Karteikarten.

\textbf{Analogie:} Wie mehrere \textbf{Bücher} gleichzeitig aufgeschlagen auf dem Tisch.

\columnbreak

\textbf{Neuen Tab öffnen:}
\begin{itemize}
\item Tippe auf das \textbf{Tabs-Symbol} (zwei Quadrate)
\item Dann auf \textbf{+} (unten)
\end{itemize}

\textbf{Zwischen Tabs wechseln:}
\begin{enumerate}
\item Tippe auf das \textbf{Tabs-Symbol}
\item Alle offenen Tabs werden angezeigt
\item Tippe auf den gewünschten Tab
\end{enumerate}

\textbf{Tab schließen:}
\begin{itemize}
\item In der Tab-Übersicht: \textbf{X} am Tab tippen
\item Oder: Tab nach \textbf{links wischen}
\end{itemize}

\vspace{3mm}

\section*{\textcolor{privategray}{Verlauf \& Privates Surfen}}

\textbf{Verlauf:} Safari merkt sich, welche Seiten du besucht hast.

\textbf{Verlauf ansehen:}\\
Buch-Symbol → Uhr-Symbol (Verlauf)

\textbf{Verlauf löschen:}\\
Einstellungen → Safari → \glqq Verlauf und Websitedaten löschen\grqq{}

\vspace{2mm}

\textbf{Privates Surfen:}
\begin{enumerate}
\item Tabs-Symbol tippen
\item Unten: Zahl tippen (z.B. \glqq 3 Tabs\grqq{})
\item \textbf{\glqq Privat\grqq{}} wählen
\end{enumerate}

$\rightarrow$ Nichts wird gespeichert (Verlauf, Cookies...)

\vspace{2mm}

\begin{center}
\fcolorbox{warnred}{warnred!8}{%
\begin{minipage}{0.9\linewidth}
\vspace{2mm}
\textbf{Achtung:}\\[2pt]
Privates Surfen macht dich nicht \textbf{unsichtbar}! Dein Internetanbieter und die Webseiten sehen dich trotzdem.
\vspace{2mm}
\end{minipage}}
\end{center}

\end{multicols}

\newpage

% ============================================================================
% SEITE 3: ÜBUNGEN & CHECKLISTE
% ============================================================================

\begin{center}
{\LARGE\bfseries Übungen \& Checkliste}\\[3pt]
{\normalsize Werde sicher im Umgang mit Safari}
\end{center}

\vspace{1mm}
\hrule
\vspace{2mm}

\begin{multicols}{2}

\textbf{\large Teil A: Begriffe zuordnen}

Was ist was? Schreibe die Nummer.

\vspace{2mm}

\renewcommand{\arraystretch}{1.5}
\begin{tabular}{@{}p{3cm}p{0.8cm}p{5cm}@{}}
1. Browser & & A. Gespeicherte Webseiten-Links \\
2. Adressleiste & & B. Mehrere Seiten gleichzeitig offen \\
3. Lesezeichen & & C. Programm zum Anzeigen von Webseiten \\
4. Tabs & & D. Feld für die Webadresse \\
5. Verlauf & & E. Liste der besuchten Seiten \\
\end{tabular}
\renewcommand{\arraystretch}{1}

\textit{\footnotesize Lösungen: 1-C, 2-D, 3-A, 4-B, 5-E}

\vspace{4mm}

\textbf{\large Teil B: Praktische Übungen}

Führe diese Aktionen in Safari aus:

\vspace{2mm}

\begin{itemize}[label=$\square$, itemsep=3pt]
\item Öffne \textbf{Safari}
\item Gib \texttt{www.tagesschau.de} in die Adressleiste ein
\item Scrolle auf der Seite nach unten
\item Tippe auf einen \textbf{Link} (z.B. einen Artikel)
\item Gehe \textbf{zurück} zur vorigen Seite
\item Speichere die Seite als \textbf{Lesezeichen}
\item Öffne einen \textbf{neuen Tab}
\item Gib \texttt{www.wetter.de} ein
\item Wechsle zwischen den \textbf{Tabs}
\item Schließe einen Tab
\end{itemize}

\vspace{4mm}

\textbf{\large Teil C: Was machst du?}

\vspace{2mm}

\textbf{Situation 1:} Du willst eine Webseite speichern, die du oft besuchst.

\rule{7cm}{0.4pt}

\vspace{3mm}

\textbf{Situation 2:} Du hast aus Versehen eine Seite geschlossen und willst sie wiederfinden.

\rule{7cm}{0.4pt}

\columnbreak

\textbf{\large Teil D: Richtig oder Falsch?}

\vspace{2mm}

\renewcommand{\arraystretch}{1.4}
\begin{tabular}{@{}p{7.5cm}cc@{}}
& R & F \\
1. Safari ist der Browser von Apple. & $\square$ & $\square$ \\
2. Man muss immer \texttt{https://} eingeben. & $\square$ & $\square$ \\
3. Lesezeichen speichern Webseiten-Adressen. & $\square$ & $\square$ \\
4. Tabs sind wie mehrere Fenster gleichzeitig. & $\square$ & $\square$ \\
5. Privates Surfen macht dich völlig anonym. & $\square$ & $\square$ \\
\end{tabular}
\renewcommand{\arraystretch}{1}

\textit{\footnotesize Lösungen: R, F, R, R, F}

\vspace{4mm}

\textbf{\large Checkliste: Das kann ich jetzt}

\vspace{2mm}

\begin{center}
\fcolorbox{safariblue}{safariblue!10}{%
\begin{minipage}{0.95\linewidth}
\vspace{2mm}
\textbf{\textcolor{safariblue}{Grundlagen}}
\begin{itemize}[label=$\square$, itemsep=2pt]
\item Ich weiß, was ein \textbf{Browser} ist
\item Ich kann eine \textbf{Webadresse eingeben}
\item Ich kann auf einer Seite \textbf{navigieren}
\item Ich kann \textbf{vor und zurück} gehen
\end{itemize}
\vspace{2mm}
\end{minipage}}
\end{center}

\vspace{2mm}

\begin{center}
\fcolorbox{bookmarkpurple}{bookmarkpurple!10}{%
\begin{minipage}{0.95\linewidth}
\vspace{2mm}
\textbf{\textcolor{bookmarkpurple}{Lesezeichen \& Tabs}}
\begin{itemize}[label=$\square$, itemsep=2pt]
\item Ich kann ein \textbf{Lesezeichen speichern}
\item Ich kann meine \textbf{Lesezeichen öffnen}
\item Ich kann \textbf{mehrere Tabs} öffnen und wechseln
\item Ich kann einen \textbf{Tab schließen}
\end{itemize}
\vspace{2mm}
\end{minipage}}
\end{center}

\vspace{2mm}

\begin{center}
\fcolorbox{bordergray}{white}{%
\begin{minipage}{0.95\linewidth}
\vspace{2mm}
\textbf{Meine wichtigsten Lesezeichen}\\[2pt]
1. \dotfill\\[2mm]
2. \dotfill\\[2mm]
3. \dotfill
\vspace{2mm}
\end{minipage}}
\end{center}

\end{multicols}

\vfill

\begin{center}
\footnotesize\textit{Stand: \today{} -- Safari ist dein Fenster ins Internet. Speichere oft besuchte Seiten als Lesezeichen!}
\end{center}

\end{document}
